\documentclass[11pt]{article}
\usepackage{fullpage, fancyhdr, ifthen, amssymb, amsfonts, amsmath, booktabs}
\usepackage{tabularx}
\usepackage{multicol}
\usepackage[margin=0.25in]{geometry}
\usepackage{upquote}
\usepackage{forest}
\usepackage{ulem}
\usepackage{colortbl}
\usepackage{enumitem}
\pagestyle{empty}

\begin{document}

\begin{table}[h]
\centering
\arrayrulecolor{black!10} % Very light table outline
\begin{tabular}{|p{\textwidth}|}
\hline
\begin{minipage}[t]{\textwidth}
\vspace{0pt}  % This is crucial for top alignment
\noindent \textbf{Instructions:} Please answer the following questions, making sure to write legibly on the following pages. The code below is a subsection of a Python file that runs Flask for a library management system.
\begin{itemize}
\item Since we are writing Python code, please be careful about spacing/indentation. Leave enough room that it is obvious when you want things indented. 
\item Make sure to write legibly.
\item All questions are worth equal points unless otherwise stated.
\item The quiz is cumulative and may include material from previous weeks.
\end{itemize}
\end{minipage} 
\\
\hline
\end{tabular}
\end{table}

{\color{lightgray}\hrule}
\begin{verbatim}
from flask import Flask, request, jsonify

app = Flask(__name__)

@app.route('/api/books/<book_id>', methods=['GET'])
def get_book(book_id):
    
    print(f"{book_id}")
    print(f"{request.args.to_dict()}")
    print(f"{request.method}")
    
    return jsonify({
        'book_id': book_id,
        'args': request.args.to_dict(),
        'method': request.method
    }), 200

@app.route('/api/books/search/<search_term>', methods=['GET'])
def search_books(search_term):
    results = get_search_results(search_term)
    # INSERT ANSWER HERE

@app.route('/api/books', methods=['POST'])
def create_book():
    data = request.get_json()
    title = data.get('title')
    author = data.get('author')
    return jsonify({'message': f'Created book: {title} by {author}'}), 201

\end{verbatim}
{\color{lightgray}\hrule}

\clearpage

\begin{tabularx}{\textwidth}{l|X}
\textbf{Quiz 6A} &   \textbf{Name: } \\
\hline
\end{tabularx}

%\vspace{1cm}

\begin{enumerate}

\item An HTTP request consists of several parts. Please identify where we would expect each of the following types of information to be located in an HTTP request:

\begin{table}[h]
\centering
\renewcommand{\arraystretch}{1.5}
\begin{tabular}{|p{7.5cm}|p{7.5cm}|}
\hline
\textbf{Information Type} & \textbf{Part of HTTP Request} \\
\hline
Search Terms from a google query & \\
\hline
Information submitted for inserting data & \\
\hline
Authentication token for API access & \\
\end{tabular}
\end{table}

\vspace{1cm}

\item Consider the \texttt{get\_book} function shown on the first page. If a client makes a GET request to the following URL:

\begin{verbatim}
GET /api/books/12345?format=json&include_reviews=true
\end{verbatim}

What would be printed to the console? Please write the exact output that would appear.

\vspace{6cm}
\end{enumerate}
\clearpage


\noindent Consider the \texttt{search\_books} function shown on the first page. The function \texttt{get\_search\_results} takes a search term and returns a list of tuples, where each tuple contains a book title and its relevance score (higher is better). 

\noindent For example, the function \texttt{get\_search\_results("monster books")} might return:

\begin{verbatim}
[
    ('Frankenstein', 0.25),
    ('Dracula', 0.75),
    ('The Strange Case of Dr. Jekyll and Mr. Hyde', 0.45), 
    ('The Island of Dr. Moreau', 0.90)
]
\end{verbatim}

\noindent The results are \textbf{NOT ordered} by relevance score, so higher relevance scores can appear before or after lower ones.

\begin{enumerate}[start=3]
\item Complete the \texttt{search\_books} function by writing code at \texttt{\# INSERT ANSWER HERE}. Your code should:

\begin{itemize}
    \item Filter the results to only include books with a relevance score greater than 0.5
    \item Create a list containing only the titles (not the scores) of the filtered books
    \item If there are results, return them using \texttt{jsonify(\{'results': <list\_of\_titles>\})} with status code 200
    \item If there are no relevant results return \texttt{jsonify(\{'results': []\})} with status code 200
\end{itemize}

\vspace{6cm}

\item Write code that finds and returns the first result from \texttt{get\_search\_results(search\_term)} that has a relevance score greater than 1.0. Your code should:

\begin{itemize}
    \item Return the {\it first} result (book title only, not the score) that has a relevance score greater than 1.0
    \item If a result is found, return them using \texttt{jsonify(\{'book': <title>\})} with status code 200
    \item If there are no results with relevance greater than 1.0, return \texttt{jsonify(\{'error': 'No results found'\})} with status code 404
\end{itemize}

\end{enumerate}

\end{document}

